% =========================================================================
% p1_sec01_introduction.tex — Part I Introduction (companion to Part II).
% Derived from the sealed monograph sec01; contributions C4/C5 routed to Part II.
% =========================================================================
\section{Introduction}
\label{sec:intro}

Insular interconnected systems occupy an extreme power engineering design space: lacking neighbor interconnections, every megawatt of reserve, reactive compensation, and damping capability must be produced within the same synchronous footprint that consumes it. The Dominican Republic is a canonical instance: a synchronous island of eleven million people without external ties, $>40\%$ of national load concentrated in the metropolitan Ozama region, an inertia anchor dominated by the Punta Catalina plant, and an instantaneous solar share reaching 36.6\% of demand in operational telemetry \cite{ocseni}. The national consequences of generation--load imbalance are evidenced by widespread service interruptions in November 2025 and February 2026 \cite{ocseni}, \cite{listin2025}, mobilizing emergency territorial response and renewing the debate on frequency stability, black-start readiness, and transmission observability \cite{blackout2004}. While the low-inertia dynamics of converter-dominated grids are theoretically established \cite{kundur1994}, \cite{tielens2016}, \cite{milano2018}, \cite{dorfler2013}, \cite{hatziargyriou2021}, empirical public instruments to analyze real insular systems remain primitive.

Power system benchmarks remain predominantly synthetic abstractions from the 1960s (WSCC 9-bus, IEEE 14-to-118-bus) with round academic reactances, fixed load triangles, and no real geography, climate, or maintenance histories. At the opposite pole, utility datasets remain encumbered by proprietary confidentiality, preventing independent re-execution. Open initiatives like MATPOWER \cite{matpower2011} and PGLib \cite{pglib2019} offer valuable instances, yet remain synthetic or anonymized. An open, complete, and measured model of an actual national transmission system with verifiable provenance represents an unmet public good for the insular Caribbean and the broader power systems community.

This article closes that gap with \textbf{SENI-142}, an open-access benchmark constructed around three commitments: (i) \textbf{Fidelity to measured reality}: a georeferenced CAD reconstruction of the Dominican bulk transmission estate validated against exactly 520 complete consecutive 168-hour weeks (87,360 hourly states spanning 3,640 Monday-to-Sunday days, 2016--2025 archive) dispatched unit-by-unit without interpolation. (ii) \textbf{Structural honesty about scale}: solving the complete national system via a hierarchical border-coordinate formulation (HBCF) in the diakoptic tradition \cite{kron1963}, \cite{happ1971}, condensed by the Schur complement \cite{zhang2005} and factorized across the ten territorial microgrids of the SNIT under Law 368-22, in harmony with General Electricity Law 125-01 \cite{ley368}. (iii) \textbf{Deterministic reproducibility}: a four-phase digital-twin lifecycle with fail-closed gates, detached SHA-512 manifests, and atomic kernel promotion guaranteeing bit-identical auditability.

The primary contributions are:
(i)~\textbf{C1 --- Benchmark:} An open, georeferenced model of the complete Dominican bulk network: 171 fenced physical utility yards, reduced to 112 transmission-level substations (345/230/138~kV) and mapped into 142 nodal computation buses (115 bulk + 27 condensed 69~kV ports), 135 transmission lines and 58 power transformers, released with topology, parameters, provenance metadata, and validation harness via GitHub.
(ii)~\textbf{C2 --- Formulation:} The HBCF diakoptic partition into ten territorial microgrids (Law 368-22) and four macro-subdomains with exact Schur condensation exhibiting complexity $O(\sum_k N_k^3 + N_B^3)$ versus monolithic $O(N^3)$ dense factorization.
(iii)~\textbf{C3 --- Deterministic digital twin:} A four-phase lifecycle (hermetic POSIX-namespace enclave with a 14-gate cascade; research-testbed static post-event $N-1$ verification against a 215-contingency catalog; SHA-512 Golden Master; and atomic kernel promotion via \texttt{renameat2} with zero regression) ensuring rigorous custodial reproducibility.
Contributions \textbf{C4} (multi-vector AC-OPF coupling hydrogen, 764~MW pumped hydro, and biomethanation) and \textbf{C5} (the 87,360-state campaign at 2.65~ms/state) are developed in the companion Part~II~\cite{seni142part2}.

The remainder of Part~I is organized as follows: Section~\ref{sec:methodology} presents the invariant formulation (graph model, 69~kV reduction, diakoptic HBCF, Schur proofs, and lifecycle); Section~\ref{sec:topology} details the 142-bus physical and microgrid topology; Section~\ref{sec:conclusion} concludes Part~I. The companion Part~II~\cite{seni142part2} develops the multi-vector AC-OPF, the 87,360-state campaign, severe contingencies, and policy implications.
